\documentclass[10pt,a4paper]{amsart}
\usepackage[margin=19mm]{geometry}
\usepackage{amsmath,amssymb,mathtools}
\usepackage[scr=rsfs,cal=euler]{mathalfa}
\usepackage{xfrac}
\usepackage[unicode,hidelinks,bookmarks=false]{hyperref}
\newcommand{\ie}{i.e.\ }
\newcommand{\cf}{cf.\ }
\newcommand{\resp}{respectively\ }
\newcommand{\Subsection}{\subsection}
\providecommand{\nobreakdash}{\nobreak\hbox{-}}
\setlength{\emergencystretch}{2em}
\multlinegap0pt
\usepackage{bm}
\usepackage{bbm}
\usepackage{dsfont}
\usepackage{xspace}
\usepackage{cancel}
\usepackage{mathtools}
\usepackage{amsthm}
\makeatletter
\@ifundefined{lemma}{\newtheorem{lemma}{Lemma}}{}
\makeatother
\title[Averages of quadratic twists of long Dirichlet polynomials]{Averages of quadratic twists of long Dirichlet polynomials}
\author{Brian Conrey, Brad Rodgers}
\date{Source: arXiv:2208.01783v2 (CC BY 4.0)}
\begin{document}
\maketitle

\noindent\emph{Source and licence.} Averages of quadratic twists of long Dirichlet polynomials. Version arXiv:2208.01783v2 (CC BY 4.0), distributed under the Creative Commons Attribution 4.0 International licence, as recorded in the arXiv licence metadata for this exact version and shown on the abstract page https://arxiv.org/abs/2208.01783v2. Full rights notice: licenses/arxiv-2208.01783-CC-BY-4.0.txt.\par\smallskip

\noindent\textit{Editorial scope.} Counted unit: the Lemma whose source label is \texttt{lemma:singularity location} in the article (source0.tex lines 682--697, source label \texttt{lemma:singularity location}), delivered together with the article's own complete original proof of it (source0.tex lines 699--725, one \texttt{proof} environment of 27 lines). No mathematics is written, repaired or re-derived: statement and proof are the article's own text line for line. Required context retained uncounted: lemma:poisson (lines 461--484); eqn:euler2 (lines 662--670). Every definition, display and label of those lines is retained unchanged. Omitted from this package and not redistributed: 1--460, 485--661, 671--681, 698--698, 726--1186. Nothing inside a retained excerpt is omitted or altered: every retained body is delivered exactly as the article's lines are written, so no omission receipt of any kind accompanies this package. Citations: the retained excerpts carry no citation command at all, so no bibliography entry of the article is transcribed and no reference is redistributed. Reference adaptation: none was needed and none was made; every reference inside the delivered unit resolves inside it, so no unresolved reference or citation is produced. Printed slips kept literal: no printed word, symbol, hypothesis or number of the counted statement or of its proof was altered. Layout and class: the article's own class is \texttt{amsart} with packages including amssymb, latexsym, amsmath, amsxtra, graphicx, subfig, xcolor; the wrapper is the lane's pinned \texttt{amsart} editorial wrapper at 10pt on A4, which restates the theorem-like environments the retained text needs and adds the article's own macro definitions that the retained text needs (none), with extra wrapper packages (bm, bbm, dsfont, xspace, cancel, mathtools). The delivered numbering is the unit's own. Assets: no retained span contains a figure environment, a graphics-inclusion command, a PDF-page inclusion, a table or a TikZ picture; no image file is copied into this package, the package declares no asset dependency and no image file is redistributed with it. Licence: version arXiv:2208.01783v2 of this article is distributed under the Creative Commons Attribution 4.0 International licence, as recorded in the arXiv licence metadata for this exact version and shown on the abstract page https://arxiv.org/abs/2208.01783v2. A full rights notice is delivered at licenses/arxiv-2208.01783-CC-BY-4.0.txt. Characters: every retained excerpt is pure ASCII, so the delivered engine, XeLaTeX with its default Latin Modern text font, reports no missing character.
	\begin{lemma}  \label{lemma:poisson} (Poisson summation formula)
		Let $m$ be a positive odd integer. We have 
		\begin{eqnarray*}
			\sum_{d ~ \textrm{odd} } \sum_{c^2\mid d\atop c\le Y} \mu(c)
			\Psi\left(\frac d D\right )\chi_d(m)
			=D\frac{\left(\frac 2 m \right)}{2m}\sum_{(c,2m)=1\atop c\le Y} \frac{\mu(c)}{c^2} \sum_{k = -\infty}^\infty (-1)^k   G_k(m) \tilde{\Psi}\left(\frac{kD}{2c^2m}\right)
		\end{eqnarray*} 
		where
		$$\tilde {\Psi}(x)=\int_{-\infty}^\infty \Psi(u)(\cos 2\pi xu+\sin 2\pi xu)~du$$
		and where $G_0(m)=\phi(m) $ if $m=\square$ and 0 otherwise; and  if $k\ne 0$ then  $G_k(m)$  is a multiplicative function of $m$ whose value at 
		$m=p^\mu$ is determined by the value of $\kappa$ for which
		$p^\kappa\mid\mid k$ by 
		\begin{eqnarray*}
			G_k(p^\mu)=\left\{
			\begin{array}{ll}
				0 & \mbox{if $ \mu \le \kappa$ and $\mu$ is odd} \\
				\phi(p^\mu) & \mbox{if $ \mu \le \kappa$ and $\mu$ is even} \\
				-p^\kappa & \mbox{if $ \mu = \kappa+1 $ and $\mu$ is even} \\
				\left(\frac{kp^{-\kappa}}{p}\right)p^\kappa\sqrt{p} & \mbox{if $ \mu = \kappa+1 $ and $\mu$ is odd} \\
				0 & \mbox{if $\mu\ge \kappa+2$}
			\end{array}
			\right.
		\end{eqnarray*}
	\end{lemma}

	\begin{align} 
		\label{eqn:euler2}
		\notag U_c(s,w) &= \sum_{(n,2c)=1} \frac{\tau_A(n)}{n^{1+w}}   
		\sum_{k=1}^\infty  (-1)^k \frac{G_{k^2}(n\ell) /\sqrt{n}}{k^{2s-2w}}\\
		\notag  &   = -(1-2^{1+2w-2s}) \sum_{(n,2c)=1} \frac{\tau_A(n)}{n^{1+w}} 
		\sum_{k=1}^\infty   \frac{G_{k^2}(n\ell) /\sqrt{n}}{k^{2s-2w}} \\
		&= - \frac{1-2^{1+2w-2s}}{1-2^{2w-2s}} \sum_{(n,2c)=1} \frac{\tau_A(n)}{n^{1+w}} 
		\sum_{k \; \mathrm{odd}}  \frac{G_{k^2}(n\ell) /\sqrt{n}}{k^{2s-2w}},
	\end{align}

	\begin{lemma}
		\label{lemma:singularity location}
		Let $a$ be a small fixed constant. Then in the region $\Re w > 0$ and $\Re(s-w) > 1/2$, for odd coprime $c$ and $\ell$,
		\begin{equation}
			\label{eq:Uc(s,w) continuation}
			U_c(s,w) = -\frac{1-2^{1+2w-2s}}{1-2^{2w-2s}} Z^{[2]}_A(1+w) \zeta^{[2]}(2s-2w) V_c(s,w;\ell),
		\end{equation}
		where $V_c(s,w;\ell)$ is analytic in $s$ and $w$ for $a \leq \Re s$ and $-a/4 \leq \Re w \leq \Re s -a/2$ and in this region (i) has a convergent Euler product over only odd primes, and (ii) satisfies $V_c(s,w;\ell) \ll_\epsilon c^\epsilon \ell^{1+\epsilon}$, where $\epsilon > 0$ can be chosen arbitrarily.
		
		Consequently for odd coprime $c$ and $\ell$, the function $U_c(s,w)$ has an analytic continuation to this region and
		\begin{equation}
			\label{eq:ResUc(s,w) continuation}
			\operatornamewithlimits{Res}_{w = -\alpha} U_c(s,w) = -\frac{1}{2} \cdot \frac{1-2^{1-2\alpha-2s}}{1-2^{-2\alpha-2s}} Z^{[2]}_{A'}(1-\alpha) \zeta^{[2]}(2s+2\alpha) V_c(s,-\alpha;\ell),
		\end{equation}
		for $A':= A \setminus \{\alpha\}$.
	\end{lemma}

	\begin{proof}
		Note that $f(k,n)=G_{k^2}(n)$ is a multiplicative function of two variables; i.e. we have $f(k_1k_2,n_1n_2)=f(k_1,n_1)f(k_2,n_2)$ whenever $(k_1n_1,k_2n_2)=1$. 
		This allows us to express $U_c(s,w)$ as an Euler product. Let
		$$
		\mathcal{H}_p(s,w;\ell) = \sum_{n,k=0}^\infty \sum_{k,n=0}^\infty \frac{\tau_A(p^n)G_{p^{2k}}(p^{n+\nu_p(\ell)})}{p^{nw+n/2+2k (s-w)}},
		$$
		where $\nu_p(\ell) := \max \{k:\; p^k | \ell\}$ is the $p$-adic valuation of $\ell$. We calculate from \eqref{eqn:euler2} that for $\Re w > 0$ and $\Re(s-w) > 1/2$, as long as $c$ and $\ell$ are coprime odd numbers,
		\begin{align}
			\label{eq:Uc(s,w) euler}
			\notag U_c(s,w) &= -\frac{1-2^{1+2w-2s}}{1-2^{2w-2s}} \prod_{p \nmid 2c\ell} \mathcal{H}_p(s,w;1) \prod_{p | c} \sum_{k=0}^\infty \frac{1}{p^{k(2s-2w)}} \prod_{p | \ell} \mathcal{H}_p(s,w;\ell) \\
			&= -\frac{1-2^{1+2w-2s}}{1-2^{2w-2s}}  \prod_{p>2} \mathcal{H}_p(s,w;1) \prod_{\substack{p|c}} \frac{ \sum_{k\geq 0} 1/p^{k(2s-2w)}}{\mathcal{H}_p(s,w;1) } \prod_{p | \ell} \frac{\mathcal{H}_p(s,w;\ell)}{\mathcal{H}_p(s,w;1)}
		\end{align}
		
		Using the values in Lemma \ref{lemma:poisson} to evaluate $G_{p^{2k}}(p^n)$, we find for $a \leq \Re s $ and $-a/4 \leq \Re w \leq \Re s - a/2$, if we let $v = 2s-2w$ (so $\Re v \geq a$),
		\begin{align*}
			\mathcal{H}_p(s,w;1) &=\Big(1 - \frac{1}{p^v}\Big)^{-1} \Big(1 + \frac{\tau_A(p^2) \phi(p^2)}{p^{3+2w} p^v} + \frac{\tau_A(p^4) \phi(p^4)}{p^{6+4w} p^{2v}} + \frac{\tau_A(p^6) \phi(p^6)}{p^{9+6w} p^{3v}}+\cdots \Big) \\
			&\quad\quad + \Big( \frac{\tau_A(p)}{p^{1+w}} + \frac{\tau_A(p^3)}{p^{2+2w}} + \frac{\tau_A(p^5)}{p^{3+5w}} +\cdots \Big)\\
			&= \Big(1 - \frac{1}{p^v}\Big)^{-1} \Big(1 + \frac{\tau_A(p)}{p^{1+w}} + \frac{\tau_A(p^2)}{p^{2(1+w)}}+\cdots\Big)\Big(1 + O\Big(\frac{1}{p^{1+\delta}}\Big)\Big),
		\end{align*}
		for sufficiently small $\delta > 0$ as long as $a$ is sufficiently small. From these terms the factor $Z^{[2]}_A(1+w) \zeta^{[2]}(2s-2w)$ can be extracted. 
		
		Note furthermore that this implies in this region of $w$ and $s$ that $\mathcal{H}_p(s,w;1) = e^{O_a(1)}$. We can use the same argument and the values in Lemma \ref{lemma:poisson} to see that $\mathcal{H}_p(s,w;\ell) \ll_a p^{\nu_p(\ell)}$. And finally it is easy to see that $\sum 1/p^{k(2s-2w)} = e^{O_a(1)}$. 
		
		If we return to \eqref{eq:Uc(s,w) euler} using this information and the fact that number of primes factors of $c$ is $o(\log c)$ (and likewise for $\ell$) we obtain \eqref{eq:Uc(s,w) continuation}.
		
		The claim \eqref{eq:ResUc(s,w) continuation} then follows immediately as $U_c(s,w)$ has a simple pole for $w$ at $-\alpha$, coming from the factor $\zeta^{[2]}(1+w+\alpha)$ in $Z^{[2]}_A(1+w)$. Note that $\zeta^{[2]}(1+w+\alpha)$ has a residue of $1/2$ at $w = -\alpha$.
	\end{proof}

\end{document}
